\documentclass[11pt]{article}
\usepackage[margin=1in]{geometry}
\usepackage{amsmath, amssymb}
\usepackage[utf8]{inputenc}
\usepackage[T1]{fontenc}
\usepackage{lmodern}
\usepackage{microtype}
\usepackage{graphicx}
\usepackage{booktabs}
\usepackage{float}
\usepackage[font=small,labelfont=bf,skip=5pt]{caption}

\title{Problem Set 9 --- Questions (m), (n) and Capstone}
\date{}

\begin{document}
\maketitle
\vspace{-3em}

% Code: ps9_part_m_n.py, ps9_capstone_invvol.py. AI log: ai_process_record.md.

%======================================================================
\section*{Part 2: MOM vs.\ TSMOM}
\subsection*{Question (m)}

We test whether one momentum strategy is redundant given the other with spanning
regressions, run in both directions for each asset class $k\in\{\text{EQ, FX, FI, CM}\}$:
\[
\text{TSMOM}^k_t=\alpha+\beta\,\text{MOM}^k_t+\gamma\,\text{RMRF}_t+\varepsilon_t,
\qquad
\text{MOM}^k_t=\alpha'+\beta'\,\text{TSMOM}^k_t+\gamma'\,\text{RMRF}_t+\varepsilon'_t .
\]
Under the null that MOM spans TSMOM, $\alpha=0$; under the null that TSMOM spans MOM,
$\alpha'=0$. We also regress diversified TSMOM on an inverse-volatility average of the
four MOM factors, and the reverse. Because each asset class has its own regressor,
GRS does not apply to the four alphas jointly, so we use a SUR-style Wald test,
$W=\hat\alpha'\hat V^{-1}\hat\alpha\sim\chi^2_4$. The sample runs from January 1985 to
May 2026 ($T=481$).

\begin{table}[H]
\centering\small
\caption{Spanning regressions, MOM vs.\ TSMOM, controlling for RMRF. Alphas in \% per year;
$t_{NW}$ uses Newey--West standard errors with 6 lags; $\rho$ is the correlation between
MOM and TSMOM. TSMOM runs at much higher volatility than MOM (e.g.\ 29\% vs.\ 4.5\% a year
in FI), so alpha magnitudes are not comparable across the two directions; the
$t$-statistics are.}
\label{tab:m}
\begin{tabular}{lrrrrrrr}
\toprule
 & \multicolumn{3}{c}{TSMOM on MOM} & \multicolumn{3}{c}{MOM on TSMOM} & \\
\cmidrule(lr){2-4}\cmidrule(lr){5-7}
 & $\alpha$ & $t$ & $t_{NW}$ & $\alpha$ & $t$ & $t_{NW}$ & $\rho$\\
\midrule
Equity indices (EQ) & 6.39 & 1.58 & 1.27 & 3.36 & 2.15 & 2.09 & 0.36\\
Currencies (FX)     & 8.46 & 4.19 & 4.20 & $-1.94$ & $-2.12$ & $-2.48$ & 0.73\\
Fixed income (FI)   & 20.04 & 4.64 & 3.99 & $-0.76$ & $-1.25$ & $-1.16$ & 0.41\\
Commodities (CM)    & 7.43 & 4.12 & 3.91 & $-3.01$ & $-1.27$ & $-1.23$ & 0.65\\
Aggregate           & 9.04 & 6.01 & 5.78 & $-1.08$ & $-1.87$ & $-1.79$ & --\\
\midrule
Joint Wald $\chi^2_4$ ($p$) & \multicolumn{3}{c}{49.29 ($<0.001$)} & \multicolumn{3}{c}{13.62 (0.009)} & \\
\bottomrule
\end{tabular}
\end{table}

\begin{itemize}
\item \textbf{Design a test to see if MOM captures TSMOM and vice versa.} We use a spanning
regression rather than a correlation, because a high correlation does not make one strategy
redundant. With one regressor, $\alpha_{B|A}>0$ exactly when $SR_B>\rho\,SR_A$, so spanning
depends on the means as well as on the correlation. The test is also asymmetric: B can have
zero alpha on A while A still has positive alpha on B, which is why we run it in both
directions.

\item \textbf{Run that test, show your results, and interpret the findings.} MOM does not
capture TSMOM: TSMOM keeps a large, significant alpha in FX, FI, CM and in aggregate
($t=6.0$), and only equities are insignificant. In the other direction, MOM's alphas on
TSMOM are zero or negative everywhere except equity index futures, where MOM still earns
about 3\% a year ($t=2.2$). The joint Wald test rejects both nulls, but the second rejection
is driven by the positive equity alpha and the negative FX alpha, so it does not mean MOM
adds value overall. TSMOM contains most of what MOM does, plus a time-series component that
MOM, being market-neutral within each asset class, cannot replicate.

\item \textbf{What do you conclude?} TSMOM largely spans MOM, but MOM does not span TSMOM.
The rule $SR_B>\rho\,SR_A$ predicts the sign of all eight alphas, because TSMOM has the
higher Sharpe ratio in every asset class and the correlations are well below one.
Cross-sectional momentum only adds something of its own in equity index futures.
\end{itemize}

%======================================================================
\section*{Part 3: Hedge Funds}
\subsection*{Question (n)}

Hedge-fund returns are in excess of the one-month T-bill rate. The spreadsheet does not
contain a T-bill column, so we use the Fama--French \texttt{RF} series. HFRI indices run
from January 1990 to August 2019 and DJCS indices from January 1994 to June 2026; each
regression ends at the last month for which all of the model's factors are available.
Dedicated Short Bias ends early (HFRI December 2017, DJCS January 2017), so each GRS test is
run on all ten indices over the shorter common sample and on the other nine over their full
sample. RMRF(lag) is the previous month's RMRF.

\begin{table}[H]
\centering
\caption{Hedge-fund alphas (\% per year) with $t$-statistics in parentheses, average
$R^2$, and GRS $F$-statistics with $p$-values. M1 CAPM; M2 adds RMRF(lag); M3 Fama--French
4 factors + lag; M4 AMP 3 factors + lag; M5 adds TSMOM; M6 adds LTREV and STR; M7 adds BAB,
Quality/Defensive and Carry.}
\label{tab:n}
\resizebox{\linewidth}{!}{%
\begin{tabular}{lrrrrrrr}
\toprule
 & M1 & M2 & M3 & M4 & M5 & M6 & M7\\
\midrule
\multicolumn{8}{l}{\emph{Panel A: HFRI}}\\
Convertible Arb   & 3.32 (3.4) & 2.58 (2.7) & 2.92 (3.0) & 2.24 (2.2) & 2.23 (2.1) & 2.23 (2.1) & 1.35 (1.4)\\
Dedicated Short   & 4.36 (2.2) & 4.54 (2.2) & 3.63 (2.2) & 0.90 (0.4) & 0.55 (0.3) & 0.47 (0.2) & $-1.08$ ($-0.5$)\\
Emerging Markets  & 3.36 (1.8) & 2.44 (1.3) & 2.28 (1.2) & 1.19 (0.6) & 1.78 (0.9) & 1.87 (0.9) & 0.05 (0.0)\\
Equity Mkt Neutral& 2.75 (5.4) & 2.70 (5.3) & 1.75 (3.9) & 1.43 (2.9) & 1.48 (2.9) & 1.51 (3.0) & 1.45 (2.9)\\
Event Driven      & 4.47 (5.7) & 3.65 (5.0) & 3.58 (5.2) & 2.83 (3.6) & 3.02 (3.8) & 3.11 (4.0) & 2.53 (3.3)\\
FI Relative Value & 4.73 (7.5) & 4.21 (6.9) & 4.27 (7.0) & 3.47 (5.4) & 3.68 (5.6) & 3.70 (5.7) & 3.23 (5.2)\\
Fund Composite    & 3.78 (5.3) & 3.23 (4.7) & 3.05 (4.9) & 2.72 (3.8) & 2.61 (3.6) & 2.70 (3.8) & 2.28 (3.2)\\
Global Macro      & 5.60 (4.6) & 5.43 (4.4) & 4.48 (3.7) & 3.93 (3.2) & 2.20 (1.8) & 2.28 (1.9) & 1.88 (1.6)\\
Long/Short Equity & 4.37 (4.6) & 3.87 (4.1) & 3.45 (4.2) & 3.31 (3.4) & 3.42 (3.5) & 3.59 (3.7) & 3.34 (3.4)\\
Managed Futures   & 4.55 (3.6) & 4.56 (3.6) & 4.05 (3.3) & 4.47 (3.4) & 2.20 (1.8) & 2.24 (1.9) & 2.40 (2.0)\\
\addlinespace
Average $|\alpha|$ & 4.13 & 3.72 & 3.35 & 2.65 & 2.32 & 2.37 & 1.96\\
Average $R^2$      & 0.38 & 0.41 & 0.49 & 0.46 & 0.48 & 0.49 & 0.53\\
GRS, all 10        & 11.02 (.000) & 10.10 (.000) & 8.40 (.000) & 6.06 (.000) & 5.02 (.000) & 4.99 (.000) & 4.82 (.000)\\
GRS, ex-Ded.\ Short & 11.24 (.000) & 10.21 (.000) & 8.68 (.000) & 6.16 (.000) & 5.12 (.000) & 5.22 (.000) & 5.00 (.000)\\
\midrule
\multicolumn{8}{l}{\emph{Panel B: DJCS}}\\
Convertible Arb   & 2.40 (2.6) & 1.52 (1.7) & 1.76 (1.9) & 1.14 (1.2) & 1.22 (1.3) & 1.24 (1.3) & 0.64 (0.7)\\
Dedicated Short   & $-0.42$ ($-0.2$) & $-0.18$ ($-0.1$) & $-0.72$ ($-0.4$) & 0.32 (0.2) & 0.44 (0.2) & 0.46 (0.2) & 0.52 (0.2)\\
Emerging Markets  & 0.93 (0.5) & $-0.19$ ($-0.1$) & $-0.33$ ($-0.2$) & $-1.71$ ($-0.9$) & $-1.50$ ($-0.8$) & $-1.41$ ($-0.7$) & $-2.03$ ($-1.1$)\\
Equity Mkt Neutral& 2.80 (4.4) & 2.69 (4.2) & 2.51 (3.8) & 1.76 (2.6) & 1.70 (2.5) & 1.71 (2.5) & 1.44 (2.1)\\
Event Driven      & 2.39 (3.0) & 1.44 (1.9) & 1.21 (1.6) & 0.18 (0.2) & 0.23 (0.3) & 0.40 (0.5) & $-0.01$ (0.0)\\
FI Relative Value & 1.75 (2.0) & 0.96 (1.1) & 0.99 (1.1) & 0.87 (1.0) & 0.94 (1.0) & 0.97 (1.1) & 0.43 (0.5)\\
Fund Composite    & 2.51 (2.9) & 1.89 (2.2) & 1.14 (1.4) & 0.42 (0.5) & $-0.08$ ($-0.1$) & 0.10 (0.1) & $-0.27$ ($-0.3$)\\
Global Macro      & 5.18 (3.5) & 5.16 (3.5) & 4.05 (2.7) & 2.61 (1.8) & 1.41 (1.0) & 1.67 (1.2) & 1.37 (1.0)\\
Long/Short Equity & 2.17 (2.2) & 1.76 (1.8) & 0.94 (1.1) & 0.44 (0.5) & 0.32 (0.3) & 0.49 (0.5) & 0.46 (0.5)\\
Managed Futures   & 3.24 (1.7) & 3.61 (1.8) & 2.27 (1.2) & 0.15 (0.1) & $-3.17$ ($-2.0$) & $-3.07$ ($-1.9$) & $-3.60$ ($-2.2$)\\
\addlinespace
Average $|\alpha|$ & 2.38 & 1.94 & 1.59 & 0.96 & 1.10 & 1.15 & 1.08\\
Average $R^2$      & 0.29 & 0.32 & 0.36 & 0.37 & 0.41 & 0.41 & 0.46\\
GRS, all 10        & 5.71 (.000) & 5.76 (.000) & 5.56 (.000) & 3.61 (.000) & 3.52 (.000) & 3.49 (.000) & 3.33 (.000)\\
GRS, ex-Ded.\ Short & 4.10 (.000) & 4.28 (.000) & 3.92 (.000) & 2.27 (.018) & 2.47 (.010) & 2.42 (.011) & 2.59 (.007)\\
\bottomrule
\end{tabular}}
\end{table}

\begin{table}[H]
\centering
\caption{Selected Model 7 loadings, with $t$-statistics in parentheses. RMRF(lag) loadings are
from Model 2.}
\label{tab:n-load}
\resizebox{\linewidth}{!}{%
\begin{tabular}{lrrrrrrr}
\toprule
 & RMRF & RMRF(lag), M2 & VAL & MOM & TSMOM & BAB & Carry\\
\midrule
\multicolumn{8}{l}{\emph{HFRI}}\\
Convertible Arb   & 0.20 (10.2) & 0.10 (5.3) & $-0.12$ ($-1.8$) & $-0.13$ ($-2.1$) & 0.00 (0.0) & 0.17 (5.6) & 0.35 (4.8)\\
Dedicated Short   & $-0.85$ ($-21.2$) & $-0.02$ ($-0.6$) & 0.84 (5.9) & 0.15 (1.1) & 0.01 (0.2) & 0.09 (1.5) & 0.25 (1.7)\\
Emerging Markets  & 0.58 (14.7) & 0.12 (3.4) & 0.02 (0.2) & 0.13 (1.0) & $-0.08$ ($-1.4$) & 0.03 (0.5) & 0.54 (3.6)\\
Equity Mkt Neutral& 0.10 (10.0) & 0.01 (0.7) & 0.04 (1.2) & 0.16 (4.8) & $-0.01$ ($-0.4$) & 0.07 (5.0) & $-0.02$ ($-0.6$)\\
Event Driven      & 0.33 (21.4) & 0.11 (7.6) & $-0.05$ ($-0.9$) & $-0.01$ ($-0.1$) & $-0.01$ ($-0.7$) & 0.09 (4.0) & 0.17 (3.0)\\
FI Relative Value & 0.16 (12.7) & 0.07 (5.9) & 0.00 (0.1) & 0.01 (0.3) & $-0.02$ ($-1.4$) & 0.10 (5.4) & 0.15 (3.3)\\
Fund Composite    & 0.36 (25.0) & 0.07 (5.4) & $-0.14$ ($-2.6$) & 0.02 (0.5) & 0.02 (1.0) & 0.04 (1.9) & 0.14 (2.6)\\
Global Macro      & 0.21 (8.8) & 0.02 (1.0) & $-0.15$ ($-1.8$) & $-0.08$ ($-1.0$) & 0.23 (7.1) & 0.07 (1.8) & 0.04 (0.4)\\
Long/Short Equity & 0.48 (24.5) & 0.07 (3.6) & $-0.17$ ($-2.5$) & 0.07 (1.1) & 0.00 ($-0.1$) & 0.04 (1.3) & 0.14 (1.9)\\
Managed Futures   & 0.22 (9.1) & 0.00 (0.0) & $-0.26$ ($-3.0$) & $-0.20$ ($-2.5$) & 0.30 (9.0) & $-0.05$ ($-1.3$) & $-0.11$ ($-1.2$)\\
\midrule
\multicolumn{8}{l}{\emph{DJCS}}\\
Convertible Arb   & 0.15 (8.1) & 0.10 (5.8) & $-0.09$ ($-1.5$) & $-0.07$ ($-1.2$) & $-0.01$ ($-0.6$) & 0.11 (4.0) & 0.27 (3.6)\\
Dedicated Short   & $-0.83$ ($-19.6$) & $-0.03$ ($-0.9$) & $-0.06$ ($-0.4$) & $-0.16$ ($-1.1$) & $-0.03$ ($-0.6$) & 0.16 (2.5) & $-0.12$ ($-0.7$)\\
Emerging Markets  & 0.43 (11.5) & 0.12 (3.7) & 0.04 (0.3) & 0.24 (1.9) & $-0.05$ ($-1.0$) & 0.00 (0.0) & 0.70 (4.6)\\
Equity Mkt Neutral& 0.10 (7.6) & 0.01 (1.0) & 0.11 (2.4) & 0.09 (2.0) & 0.01 (0.7) & 0.00 ($-0.2$) & 0.08 (1.4)\\
Event Driven      & 0.30 (20.6) & 0.11 (7.6) & 0.04 (0.9) & 0.08 (1.6) & $-0.01$ ($-0.5$) & 0.08 (3.4) & 0.17 (2.9)\\
FI Relative Value & 0.11 (6.4) & 0.09 (5.6) & $-0.15$ ($-2.7$) & $-0.11$ ($-2.1$) & $-0.02$ ($-0.7$) & 0.13 (5.0) & 0.29 (4.2)\\
Fund Composite    & 0.29 (18.9) & 0.07 (4.3) & $-0.13$ ($-2.5$) & 0.07 (1.3) & 0.09 (4.0) & 0.08 (3.1) & 0.23 (3.5)\\
Global Macro      & 0.22 (8.0) & 0.00 (0.1) & $-0.05$ ($-0.6$) & 0.00 (0.0) & 0.22 (5.7) & 0.14 (3.2) & 0.24 (2.1)\\
Long/Short Equity & 0.45 (25.4) & 0.05 (2.5) & $-0.17$ ($-2.9$) & 0.23 (3.9) & 0.01 (0.5) & 0.03 (1.1) & 0.21 (2.9)\\
Managed Futures   & 0.06 (1.9) & $-0.04$ ($-1.1$) & 0.01 (0.1) & $-0.24$ ($-2.3$) & 0.61 (13.8) & 0.05 (0.9) & 0.00 (0.0)\\
\bottomrule
\end{tabular}}
\end{table}

\begin{itemize}
\item \textbf{Which model performs the best at pricing the hedge fund portfolios and why?}
Model 7 performs best, although every model is rejected by the GRS test. It has the
smallest average alpha for HFRI and the highest average $R^2$ in both families, because it
adds the factors that credit- and carry-style strategies actually trade. Most of the
improvement comes earlier, though: moving to the AMP factors (M4) and adding TSMOM (M5)
roughly halves the alphas of Global Macro and Managed Futures, while the reversal factors
in M6 add almost nothing. The everywhere factors and TSMOM do the heavy lifting, and BAB,
Quality and Carry finish the job.

\item \textbf{What factors do the various hedge fund strategies load on?} Directional
strategies load mainly on the market: Long/Short Equity and Emerging Markets have betas of
about 0.45--0.6, and Dedicated Short about $-0.85$. Managed Futures and Global Macro load on
TSMOM, Convertible Arbitrage, FI Relative Value and Event Driven load on BAB and Carry, and
Equity Market Neutral loads on momentum. Each strategy looks like a combination of a few
traded factors.

\item \textbf{Do these make sense?} Yes, the loadings match what the strategies say they do.
Managed futures funds run trend-following programs, which is time-series momentum;
convertible and fixed-income arbitrage earn a credit and carry spread, which shows up as
Carry and BAB; and a short-biased fund should have a strongly negative market beta. The
factor models recover the strategies' stated styles from returns alone.

\item \textbf{Is there anything surprising?} Yes, three things stand out. Long/Short
Equity and the fund composites load \emph{negatively} on VAL, so these funds tilt toward
growth rather than value. Equity Market Neutral is not market neutral (beta about 0.1) and keeps a
significant alpha under every model. Finally, DJCS Managed Futures has a \emph{negative}
alpha of about $-3.6\%$ a year under Model 7 ($t=-2.2$): it earns less than its TSMOM
exposure would deliver, which is what we would expect from fees charged on a cheap,
replicable strategy.

\item \textbf{What does the loading on the lagged market return tell you?} It signals stale
prices in illiquid holdings. The lagged loading is significant (about 0.06--0.12) for
Convertible Arbitrage, FI Relative Value, Event Driven, Emerging Markets and the fund
composites, whose assets are marked infrequently, so their reported returns respond to the
market with a delay. Their true beta is closer to the sum of the two loadings, which means
the CAPM understates their market risk and overstates their alpha. The lagged loading is
near zero for Global Macro and Managed Futures, which trade liquid futures and are marked
to market every day.
\end{itemize}

%======================================================================
\section*{Part 4: Capstone --- Build and Stress-Test a Multi-Factor Combination}

\subsection*{Steps 1--2: Rule and hypothesis (written before running)}

\begin{itemize}
\item \textbf{Which method we chose and why.} We chose inverse-volatility weighting,
$w_i\propto 1/\hat\sigma_i$, over a ten-factor menu: RMRF, SMB, VAL$^{\text{everywhere}}$,
MOM$^{\text{everywhere}}$, TSMOM, BAB, Quality/Defensive, Carry, STR and LTREV. The rule uses
only volatilities, which are estimated far more precisely than means, so it should capture
most of the diversification benefit with little estimation error. The benchmarks are equal
weight and the in-sample tangency portfolio, $w\propto\hat\Sigma^{-1}\hat\mu$.

\item \textbf{Predicted Sharpe ratios.} We predicted an in-sample Sharpe ratio of about 1.1
and an out-of-sample Sharpe ratio of about 1.05 for inverse volatility (gap about 0.05),
about 1.0 with no gap for equal weight, and about 1.9 in sample against 1.4 out of sample
for tangency (gap about 0.5).

\item \textbf{Which method we expected to degrade most.} We expected tangency to degrade
most. With ten correlated factors it leans heavily on noisy mean estimates and takes large
offsetting positions in correlated pairs such as VAL and LTREV, or MOM and TSMOM.

\item \textbf{Risk metric beyond Sharpe.} We chose maximum drawdown. Inverse volatility
overweights low-volatility factors such as Quality and Carry, which tend to crash together
in liquidity events, so we expected its drawdown to look worse than its Sharpe ratio
suggests.
\end{itemize}

\subsection*{Step 3: Implementation}

The common sample is January 1985 to February 2026 ($T=494$), limited by TSMOM at the start
and by Quality and Carry at the end. Every portfolio is rescaled to 10\% annual volatility,
so that drawdowns are comparable and not driven by leverage. For the out-of-sample version,
weights estimated on even months are applied to odd months and vice versa, and each
held-out half is scaled using the volatility of its estimation half only. This is a
cross-fitted, held-out estimate, not a real-time backtest, because the estimation half
contains months that come later in calendar time than the month being scored. The AMP
spanning regression is $r_{p,t}=\alpha+\beta_1\text{RMRF}_t+\beta_2\text{VAL}_t+\beta_3\text{MOM}_t+\varepsilon_t$.
The LLM-generated code reproduced our estimates exactly; the corrections we made are listed
in the AI workflow log.

\subsection*{Step 4: Results}

\begin{table}[H]
\centering\small
\caption{Capstone results, all portfolios scaled to 10\% annual volatility. Sharpe ratios are
annualized. OOS uses the even/odd-month cross-fit. Maximum drawdown is for the OOS series,
with peak, trough and recovery months. AMP alpha in \% per year with Newey--West (6 lags)
$t$-statistics in parentheses; OLS $t$-statistics are 2.33, 1.98, 5.29 (IS) and 2.23, 1.97,
1.89 (OOS).}
\label{tab:cap}
\resizebox{\linewidth}{!}{%
\begin{tabular}{lrrrrlrr}
\toprule
 & IS Sharpe & OOS Sharpe & Gap & OOS MDD & MDD path & AMP $\alpha$, IS & AMP $\alpha$, OOS\\
\midrule
Inverse volatility & 1.28 & 1.27 & 0.01 & $-23.0\%$ & 2018-09 $\to$ 2020-03 $\to$ 2021-05 & 2.79 (1.91) & 2.67 (1.83)\\
Equal weight       & 1.22 & 1.22 & 0.00 & $-26.3\%$ & 2007-06 $\to$ 2009-05 $\to$ 2010-12 & 2.24 (1.67) & 2.23 (1.66)\\
Tangency           & 1.68 & 1.31 & 0.37 & $-18.1\%$ & 1987-09 $\to$ 1987-12 $\to$ 1988-07 & 4.77 (4.45) & 1.97 (1.68)\\
\bottomrule
\end{tabular}}
\end{table}

\begin{itemize}
\item \textbf{How large was the out-of-sample degradation for your rule versus the
benchmarks? Did the method you predicted would degrade most actually degrade most?} Our rule
barely degraded, and tangency degraded most, as we predicted. Inverse volatility and equal
weight only need volatilities, which are almost identical across the two halves, so their
Sharpe ratios hardly move. Tangency loses about 0.37 of Sharpe ratio, and its weights are
unstable: Carry has a weight of $+0.27$ when estimated on even months and $-0.10$ on odd
months. Even after degrading, tangency still edges out inverse volatility out of sample, so
estimation error shrank its advantage rather than reversing it. Our predicted gaps were close,
but we underestimated the level of every Sharpe ratio.

\item \textbf{Did your rule earn a significant alpha over the AMP 3-factor model, or does
the everywhere combination already span it? What does this say about whether your added
complexity was worth it?} No, the alpha of about 2.7\% a year is economically meaningful but
not robustly significant ($t=1.83$ with Newey--West errors). The AMP model explains about
half of its variance, and the remaining return comes from factors outside the model, mainly
Quality, Carry and TSMOM. Tangency's large in-sample alpha disappears once its weights are
estimated on the other half of the sample. Extra complexity therefore buys nothing out of
sample: the simple volatility-based rules do as well as the optimizer, and none of them
clearly beats the everywhere factors.

\item \textbf{On your chosen non-Sharpe risk metric, did your rule look worse than its
Sharpe ratio suggested? Connect this to the non-normality you documented in Part 1.} Yes,
the drawdown is worse than a Sharpe ratio of 1.27 suggests. At 10\% volatility the
portfolio lost 23\% from September 2018 to March 2020 and needed until May 2021 to recover,
almost three years in all. This matches the non-normality from Part 1: the combination
inherits the factors' negative skewness (about $-1.0$) and fat tails (excess kurtosis about
7), and diversification across factors does not remove crash risk when the factors fall
together in a liquidity event.

\item \textbf{What is the single most important limitation of your backtest?} Data snooping
in the choice of factor menu is the most important limitation. All ten factors are
published, and they were published because they worked over a sample that overlaps ours,
so even the out-of-sample Sharpe ratios are optimistic. The even/odd split also assumes the
factor returns are stationary over 40 years and is not a real-time test, and we ignore
transaction costs, which fall hardest on STR and the high-turnover momentum factors.
\end{itemize}

\end{document}
